	\subsection{7. Discussion}\label{discussion}}

\hypertarget{interpretation-of-results}{%
	\subsubsection{7.1 Interpretation of
		Results}\label{interpretation-of-results}}

A2C outperforms Double DQN on moving IoT service composition while
preserving STR-based selection. Three factors drive improvements.

First, actor-critic gives stable learning by combining value estimation
with direct policy optimization. The advantage function reduces variance
while keeping updates unbiased, so the agent learns effectively from
fewer samples.

Second, proactive composition via trajectory prediction anticipates
service positions instead of just reacting. The A2C agent picks services
that maintain capacity across the composition horizon, not just
currently available ones.

Third, separate networks with LSTM trajectory encoding specialize
processing for spatio-temporal states. More parameters and training
time, but better performance when movement patterns matter.

The STR-based selection continues driving decisions---the agent learns
to maximize expected capacity while maintaining stability.

\hypertarget{comparison-with-prior-work}{%
	\subsubsection{7.2 Comparison with Prior
		Work}\label{comparison-with-prior-work}}

A2C improves over Double DQN on both datasets.

Random Waypoint: A2C Separate achieves 95.2\% versus 92.4\% at low
mobility (3\% improvement). At high mobility (10 km/h), the gap grows to
13.9\% (85.7\% vs 71.8\%).

Vehicle Routes: At 80 km/h, A2C Separate reaches 73.5\% versus 54.3\%
for Double DQN---a 19.2\% absolute improvement. Structured movement
patterns help trajectory prediction.

Capacity satisfaction improves because A2C anticipates degradation and
picks services with better margin.

\hypertarget{practical-implications}{%
	\subsubsection{7.3 Practical
		Implications}\label{practical-implications}}

\textbf{Edge Deployment}: A2C with shared networks balances performance
and computation for edge deployment. Inference cost allows real-time
decisions on edge devices.

\textbf{Stability}: A2C re-composes 69.2/hr versus 124.3/hr for Double
DQN on waypoint data. Less disruption, lower overhead for production
systems.

\textbf{STR Quality}: STR-derived capacity provides a grounded service
quality metric tied to spatial proximity.

\hypertarget{limitations}{%
	\subsubsection{7.4 Limitations}\label{limitations}}

Three limitations point to future work:

\textbf{STR Model}: Uses simplified distance-based models. More complex
propagation---multipath fading, shadowing, interference---needs
investigation.

\textbf{Centralized}: Assumes centralized composition. Distributed
multi-agent approaches could scale better for large IoT systems.

\textbf{Validation}: We tested on real GPS from Illinois, but AP
locations are simulated. Future work should use real IoT service
availability data.

\begin{center}\rule{0.5\linewidth}{0.5pt}\end{center}

\hypertarget{conclusion}{%
